\documentclass[../../main.tex]{subfiles}

\graphicspath{{../../image/}{../../figure/}}

\begin{document}

\section{条件独立性}\label{section:ch11-条件独立性}

多变量概率分布的一个重要概念是条件独立性（Dawid, 1979）。考虑3个变量$a$、$b$、$c$，假设在给定$b$和$c$的条件下，$a$的条件分布不依赖于$b$的值，即

\begin{align}
  p(a\mid b,c) = p(a\mid c)
  \label{align:ch11-条件独立性定义}
\end{align}

我们称之为，在给定$c$的条件下，$a$与$b$是条件独立的。如果我们考虑以$c$为条件的$a$和$b$的联合分布，就可以用稍微不同的方式来表示，形式如下：

\begin{align}
  p(a,b\mid c) = p(a\mid b,c)\,p(b\mid c) = p(a\mid c)\,p(b\mid c)
  \label{align:ch11-条件独立性分解}
\end{align}

其中利用了概率的乘积法则和式~\eqref{align:ch11-条件独立性定义}。我们看到，以$c$为条件的$a$和$b$的联合分布可以分解为$a$的边缘分布和$b$的边缘分布的乘积（后者同样都以$c$为条件）。也就是说，在给定$c$的条件下，变量$a$和$b$在统计上是独立的。请注意，我们对条件独立性的定义要求式~\eqref{align:ch11-条件独立性定义}或式~\eqref{align:ch11-条件独立性分解}必须对$c$的每一个可能值都成立，而不仅仅对$c$的某些值成立。我们有时也会使用条件独立性的速记符号（Dawid, 1979），其中

\begin{align}
  a \perp b \mid c
  \label{align:ch11-条件独立性记号}
\end{align}

表示在给定$c$的条件下，$a$和$b$是条件独立的。条件独立性在机器学习的概率模型中发挥着重要作用，因为它既简化了模型的结构，也简化了在模型下执行推理和学习所需的计算。

如果能用条件分布的乘积（即有向图的数学表达式）得出一组变量的联合分布表达式，则原则上可以通过反复应用概率的加和法则与乘积法则来检验任何潜在的条件独立性是否成立。实际上，该方法十分耗时。图模型的一个重要而优雅的特点是，可以直接从图中读取联合分布的条件独立性，而无须执行任何分析操作。实现此过程的一般框架称为d分离（d-separation），其中“d”代表“有向”（Pearl, 1988）。本节将介绍d分离的概念并给出d分离准则的一般说明，形式证法见 Lauritzen (1996)。

\subsection{3个示例图}\label{subsection:ch11-3个示例图}

在开始讨论有向图的条件独立性之前，我们先来看3个简单的例子，每个例子都涉及只有3个节点的图。这些将共同促进和说明d分离的关键概念。第1个示例如图11.14所示，利用一般结果[式~\eqref{align:ch11-有向图联合分布}]可以很容易地写出该图所对应的联合分布，即

\begin{align}
  p(a,b\mid c) = p(a\mid c)\,p(b\mid c)
  \label{align:ch11-第一示例联合}
\end{align}

如果未观测到任何变量，则可以通过将式~\eqref{align:ch11-第一示例联合}的两边相对于$c$边缘化来研究$a$和$b$是否独立，即

\begin{align}
  p(a,b) = \sum_c p(a\mid c)\,p(b\mid c)\,p(c)
  \label{align:ch11-第一示例边缘化}
\end{align}

一般来说，这不会分解为乘积$p(a)p(b)$，因此

\begin{align}
  a \not\perp b \mid \emptyset
  \label{align:ch11-第一示例非条件独立}
\end{align}

其中$\emptyset$表示空集，符号上表示条件独立性一般不成立。当然，对于特定的分布来说，条件独立性也可能由于与各种条件概率相关的特定数值而成立，但一般来说，这并不是由图的结构得出的。

假设我们以节点$c$为条件，如图11.15中的有向图所示。根据式~\eqref{align:ch11-第一示例联合}，我们可以很容易地写出在给定$c$的条件下$a$和$b$的条件分布，其形式为

\begin{align}
  p(a,b\mid c) = p(a\mid c)\,p(b\mid c)
\end{align}

于是，我们得到如下条件独立性：

\begin{align}
  a \perp b \mid c
\end{align}

可以考虑从节点$a$经过节点$c$到达节点$b$的路径，从而对这一结果进行简单的图解释。节点$c$之所以相对于该路径称为尾尾相接（tail-to-tail）节点，就是因为该节点与两个箭头的尾部相连，而连接节点$a$和$b$的路径的存在会导致这两个节点相互依赖。然而，如图11.15所示，当我们以节点$c$为条件时，条件节点会“阻塞”从$a$到$b$的路径，并使$a$和$b$变得（条件）独立。

同样，我们也可以考虑图11.16所示的有向图。根据式~\eqref{align:ch11-有向图联合分布}，我们可以得到与该有向图相对应的联合分布，即

\begin{align}
  p(a,b,c) = p(a)\,p(c\mid a)\,p(b\mid c)
  \label{align:ch11-第二示例联合}
\end{align}

假设未观测到任何变量。我们可以通过对$c$进行边缘化来检验$a$和$b$是否独立，

\begin{align}
  p(a,b) = p(a)\sum_c p(c\mid a)\,p(b\mid c) = p(a)\,p(b\mid a)
\end{align}

从而得出

\begin{align}
  a \not\perp b \mid \emptyset
\end{align}

假设以节点$c$为条件，如图11.17中的有向图所示。利用贝叶斯定理和式~\eqref{align:ch11-第二示例联合}，即可得出

\begin{align}
  p(a,b\mid c) = \frac{p(a)\,p(c\mid a)\,p(b\mid c)}{p(c)} = p(a\mid c)\,p(b\mid c)
\end{align}

于是，我们再次得到如下条件独立性：

\begin{align}
  a \perp b \mid c
\end{align}

和以前一样，我们可以用图形来解释这些结果。就节点$a$到节点$b$的路径而言，节点$c$称为头尾相接（head-to-tail）节点。该路径连接了节点$a$和$b$，并使它们相互依赖。假设我们现在观测到了节点$c$，如图11.17所示，那么这个观测结果就会“阻塞”从节点$a$到$b$的路径，从而得到条件独立性$a\perp b\mid c$。

最后，让我们来看看第3个示例，如图11.18中的有向图所示。利用一般结果[式~\eqref{align:ch11-有向图联合分布}]写出该图所对应的联合分布，即

\begin{align}
  p(a,b,c) = p(a)\,p(b)\,p(c\mid a,b)
  \label{align:ch11-第三示例联合}
\end{align}

首先考虑未观测到任何变量的情况。在$c$上对式~\eqref{align:ch11-第三示例联合}的两边进行边缘化，得到

\begin{align}
  p(a,b) = p(a)\,p(b)
\end{align}

因此，与前两个例子所不同的是，$a$和$b$是独立的，且未观测到任何变量。我们可以将这一结果写成

\begin{align}
  a \perp b \mid \emptyset
\end{align}

假设我们以节点$c$为条件，如图11.19所示，则$a$和$b$的条件分布如下：

\begin{align}
  p(a,b\mid c) = \frac{p(a)\,p(b)\,p(c\mid a,b)}{p(c)}
\end{align}

一般来说，这不会分解为乘积$p(a\mid c)\,p(b\mid c)$，因此

\begin{align}
  a \not\perp b \mid c
\end{align}

第3个示例的行为与前两个示例相反。观察图11.18，节点$c$相对于从节点$a$到节点$b$的路径是头头相接（head-to-head）的。节点$c$有时也称碰撞节点。当未观测到节点$c$时，就会“阻塞”路径，且导致$a$和$b$是独立的。然而，以$c$为条件会“疏通”路径，并使$a$和$b$相互依赖。

\subsection{相消解释}\label{subsection:ch11-相消解释}

我们有必要花点时间进一步了解图11.19中的有向图的异常行为。如图11.20所示，考虑这样一个图，该图对应的问题与汽车燃油系统有关的3个二进制随机变量。其中，变量$B$代表电池的状态，即满电（$B=1$）或没电（$B=0$）；变量$F$代表油箱的状态，即满油（$F=1$）或没油（$F=0$）；变量$G$代表电动燃油表的读数，$G=1$表示油箱满油，$G=0$表示油箱没油，电池要么满电要么没电，油箱要么满油要么没油，且先验概率为

\begin{align}
  p(F=1)=0.9,\qquad p(B=1)=0.9
\end{align}

在给定油箱状态和电池状态的情况下，电动燃油表读数为1的概率为

\begin{align}
  p(G=1\mid B=0,F=0)=0.1,\quad
  p(G=1\mid B=0,F=1)=0.2,\quad
  p(G=1\mid B=1,F=0)=0.2,\quad
  p(G=1\mid B=1,F=1)=0.8
\end{align}

可见这是一个相当不靠谱的燃油表！因为剩下的所有概率都由概率和为1这一要求决定。

在我们观测到任何数据之前，油箱为空的先验概率为$p(F=0)=0.1$。假设我们观测燃油表并发现它的读数为零，即$G=0$，与图11.20的中间图相对应。我们可以使用贝叶斯定理来计算油箱为空的后验概率。首先计算

\begin{align}
  p(G=0) = \sum_B\sum_F p(G=0\mid B,F)\,p(B)\,p(F) = 0.315
  \label{align:ch11-燃油表边缘}
\end{align}

然后计算

\begin{align}
  p(F=0\mid G=0) = \frac{p(G=0\mid F=0)\,p(F=0)}{p(G=0)} \approx 0.257
  \label{align:ch11-油箱空后验1}
\end{align}

所以$p(F=0\mid G=0)>p(F=0)$，从而只要测到燃油表读数为零，就更有可能确定油箱确实是空的，正如我们潜意识里所认为的那样。接下来，假设我们还检查了电池的状态并发现电池没电，即$B=0$。如图11.20的右图所示，我们现在已经观测到燃油表和电池的状态。根据对燃油表状态和电池状态的观测，油箱为空的后验概率为

\begin{align}
  p(F=0\mid G=0,B=0) = \frac{p(G=0\mid B=0,F=0)\,p(F=0)}{p(G=0\mid B=0)} \approx 0.111
  \label{align:ch11-油箱空后验2}
\end{align}

其中，先验概率$p(B=0)$在式~\eqref{align:ch11-油箱空后验2}的分子和分母之间相互抵消。因此，由于观测到电池的状态，油箱为空的概率降低了（从0.257降至0.111）。这与我们的直觉相吻合，即电池没电可以相消解释（explain away）燃油表读数为零的观测结果。通过观测燃油表的读数，我们发现油箱和电池的状态确实是相互依赖的。事实上，如果我们不直接观测燃油表，而是观测节点$G$的某个后代节点的状态，例如一个相当不可靠的目击者说他看到燃油表的读数为零，则情况也会如此（见习题11.14）。请注意，概率$p(F=0\mid G=0,B=0)\approx 0.111$大于先验概率$p(F=0)=0.1$，因为观测到燃油表读数为零仍然提供了一些支持油箱为空的证据。

总而言之，在未观测到尾对尾节点或头对尾节点的情况下，路径会变得畅通；而一旦尾对尾节点或头对尾节点被观测到，路径就会变得阻塞。相比之下，头对头节点在未被观测到时会阻塞路径，而一旦头对头该节点和/或其至少一个后代被观测到，路径就会变得畅通。

第3个例子还有一个微妙之处。首先，我们来介绍一些术语。如果存在一条从$x$到$y$的路径，且这条路径上的每一步都遵循箭头的方向，则可以说节点$y$是节点$x$的后代。然后可以证明，如果节点或其任何后代被观测到，则头头相接的路径就会变得畅通（见习题11.13）。

\subsection{d分离}\label{subsection:ch11-d分离}

本小节将对有向图的d分离特性（Pearl, 1988）进行一般性阐述。考虑一个有向图，其中$\mathbf{A}$、$\mathbf{B}$、$\mathbf{C}$是任意不相交的节点集（但它们的并集可能小于图中的完整节点集）。我们希望确定给定的有向无环图是否隐含特定的条件独立性$\mathbf{A}\perp\mathbf{B}\mid\mathbf{C}$。为此，我们需要考虑从$\mathbf{A}$中任何节点到$\mathbf{B}$中任何节点的所有可能路径。如果任何这样的路径包含一个节点，则任何这样的路径可称为被阻塞路径，并且可能出现如下两种情况之一：

\begin{enumerate}
  \item 路径上的箭头在节点处头尾相接或尾尾相接，且节点位于集合$\mathbf{C}$中；
  \item 箭头在节点处头头相接，且节点及其任何后代节点都不在集合$\mathbf{C}$中。
\end{enumerate}

如果所有路径都被阻塞，那么$\mathbf{A}$和$\mathbf{B}$就会被$\mathbf{C}$分离，图中所有变量的联合分布将满足$\mathbf{A}\perp\mathbf{B}\mid\mathbf{C}$。

d分离法如图11.21所示。在图11.21(a)中，从$a$到$b$的路径既没有被节点$f$阻塞，因为节点$f$是这条路径的尾对尾节点，且未被观测到；也没有被节点$e$阻塞，因为尽管节点$e$是头对头节点，但它在条件集中只有一个后代节点$c$。因此，条件独立性$\mathbf{A}\perp\mathbf{B}\mid\mathbf{C}$从该图中无法得出。在图11.21(b)中，从$a$到$b$的路径被节点$f$阻塞，因为这是一个被观测到的尾对尾节点，所以根据该图进行分解的任何分布都将满足条件独立性$a\perp b\mid f$。请注意，这条路径也被节点$e$阻塞了，因为节点$e$是一个头对头节点，节点$e$及其子节点都不在条件集合中。在d分离中，图11.12中以浮动变量表示的$\lambda$等这些节点在d分离中不起作用。参数的行为与观测节点的行为相同。然而，这些节点不存在边缘分布，因此参数节点本身永远不会有父节点，经过这些节点的所有路径永远是尾尾相接的，因此会被阻塞。

独立同分布数据是条件独立性和d分离的另一个示例（参见2.3.2小节）。观察图11.12所示的二元分类模型。这里的随机节点分别对应$t_n$、$\mathbf{w}$和$t$。我们可以看到，对于从子到任何一个节点$t_n$的路径来说，$\mathbf{w}$的节点是尾对尾的，因此我们具有以下条件独立性：

\begin{align}
  t \perp \{t_n\} \mid \mathbf{w}
  \label{align:ch11-独立同分布条件独立}
\end{align}

因此，以网络参数$\mathbf{w}$为条件，$\hat{t}$的预测分布与训练数据$\{t_1,\cdots,t_N,\mathbf{x}_n\}$无关。我们可以首先使用训练数据来确定$\mathbf{w}$的后验分布（或后验分布的近似值），然后丢弃训练数据，并使用$\mathbf{w}$的后验分布来预测新输入的观测结果中的$t$。

\subsection{朴素贝叶斯}\label{subsection:ch11-朴素贝叶斯}

有一种相关的图结构出现在朴素贝叶斯（Naive Bayes）模型中，其中我们使用条件独立性假设来简化模型结构。假设我们的数据由向量$\mathbf{x}$的观测值组成，我们希望将$\mathbf{x}$的观测值归入$K$个类别中的一个类别（后文简称类）。我们可以为每个类定义类-条件密度$p(\mathbf{x}\mid C_k)$以及先验类概率$p(C_k)$。朴素贝叶斯模型的关键假设是，以类$C_k$为条件，将输入变量的分布分解为两个或多个密度的乘积。假设我们把$\mathbf{x}$分成$L$个元素，即$\mathbf{x}=(x_1,\cdots,x_L)$，则朴素贝叶斯模型的形式为

\begin{align}
  p(\mathbf{x}\mid C_k) = \prod_{j=1}^{L} p(x_j\mid C_k)
  \label{align:ch11-朴素贝叶斯分解}
\end{align}

假设式~\eqref{align:ch11-朴素贝叶斯分解}分别对每个类$C_k$成立。图11.22是该模型的图形化表示。我们可以看到，在$j\neq i$的情况下，给定$C$时，$x_i$和$x_j$是条件独立的。因为这种路径在节点$C$处是尾尾相接的，所以当给定$C$时，对$C$的观测会阻塞$x_i$和$x_j$之间的路径。然而，如果我们将$C$去边缘化，从$x_i$到$x_j$的尾尾相接路径将不再受阻。所以一般来说，边缘密度$p(\mathbf{x})$不会相对于元素$(x_1,\cdots,x_L)$进行分解。

如果我们得到一个标记训练集，其中包括观测值$\{\mathbf{x}_1,\cdots,\mathbf{x}_N\}$以及它们的类标签，我们就可以假设数据是从模型中独立提取的，从而使用最大似然法将朴素贝叶斯模型拟合到训练数据中（见习题11.15）。通过使用相应的标记数据分别拟合每个类的模型，然后将先验类概率$p(C_k)$设置为等于每个类中训练数据点的比例，从而得到解决方案。向量$\mathbf{x}$属于类$C_k$的概率由贝叶斯定理给出，其形式为

\begin{align}
  p(C_k\mid \mathbf{x}) = \frac{p(\mathbf{x}\mid C_k)\,p(C_k)}{p(\mathbf{x})}
  \label{align:ch11-朴素贝叶斯后验}
\end{align}

其中$p(\mathbf{x}\mid C_k)$由式~\eqref{align:ch11-朴素贝叶斯分解}给出，$p(\mathbf{x})$可以用下式求出：

\begin{align}
  p(\mathbf{x}) = \sum_k p(\mathbf{x}\mid C_k)\,p(C_k)
  \label{align:ch11-朴素贝叶斯边缘}
\end{align}

用于二维数据空间的朴素贝叶斯模型如图11.23所示，其中$\mathbf{x}=(x_1,x_2)$。此处，我们假设两个类中每个类的类-条件密度$p(\mathbf{x}\mid C_k)$都是轴对齐的高斯密度，因此它们各自将对$x_1$和$x_2$进行分解，从而

\begin{align}
  p(\mathbf{x}\mid C_k) = p(x_1\mid C_k)\,p(x_2\mid C_k)
  \label{align:ch11-二维朴素贝叶斯}
\end{align}

然而，边缘密度$p(\mathbf{x})$由下式给出：

\begin{align}
  p(\mathbf{x}) = \sum_k p(\mathbf{x}\mid C_k)\,p(C_k)
  \label{align:ch11-二维朴素贝叶斯边缘}
\end{align}

它现在是高斯混合模型且不会对$x_1$和$x_2$进行分解。在融合不同来源的数据（如用于医疗诊断的血液检测和皮肤图像）时，将会涉及朴素贝叶斯模型的简单应用（参见5.2.4小节）。当输入空间的维度$D$较高时，在全$D$维空间中进行密度估计将更具挑战性，朴素贝叶斯假设会有所帮助。如果输入向量同时包含离散变量和连续变量，这种方法也会很有用，因为每种变量都可以用适当的模型来表示（例如，对于二元观测可以使用伯努利分布，对于实值变量可以使用高斯分布）。该模型的条件独立性假设显然是强假设，可能会导致对类-条件密度的表示相当糟糕。然而，即使这一假设没有得到精确的满足，模型在实际应用中也仍然可以提供良好的分类性能，因为决策边界可能对类-条件密度中的某些细节不敏感。

\subsection{生成式模型}\label{subsection:ch11-生成式模型}

机器学习的许多应用都可以看作逆问题（inverse problem）的示例，其中有一个产生数据的底层过程（通常是物理过程），我们现在的目标是学习如何反转这个过程。一种方法是训练深度神经网络，如卷积网络，从而将图像作为输入并生成描述对象的类型、位置和大小的输出。因此，作为判别模型的一个示例，这种方法试图直接解决逆问题。只要有足够多的标记图像，这种方法就能达到很高的准确率。在实践中，未标记的图像往往很多，而在获得训练集的过程中，大部分工作量耗费在证明标记上，这个证明过程可能需要手动完成。我们的判别模型在训练过程中无法直接使用未标记的图像。

另一种方法是对生成过程进行建模，然后进行计算反转。在我们的图像示例中，如果假设对象的类、位置和大小都是独立选择的，则可以使用有向图模型来表示生成过程，如图11.24所示。请注意，箭头的方向与生成步骤的顺序相对应，因此该模型代表了观测数据生成的因果过程（Pearl, 1988）。这是生成式模型的一个示例，因为一旦经过训练，它就可以用来生成合成图像，方法是首先从学习到的先验分布中选择对象的类、位置和大小，然后从学习到的条件分布中抽取图像样本。对象所属的类（离散变量）与对象的位置和方向（连续变量）具有独立的先验概率。图像（像素强度阵列）的概率分布取决于对象所属的类及其位置和方向。

根据图11.24中的图模型，在未观测到图像时，类、位置和大小变量是独立的。这是因为这些变量中的任何两个变量之间的每条路径就图像变量而言是头头相接的，这是未观测到的。然而，当我们观测到一幅图像时，这些路径就会变得畅通，类、位置和大小变量也不再独立。从直觉上讲，这是合理的，因为被告知图像中对象所属的类可以为我们提供相关信息，从而帮助我们确定对象的位置和大小。例如，一个对象的图像可以看作一个生成过程的输出，在这个生成过程中，对象的类型、位置和方向也从一些先验分布中选择，然后创建图像。在给定一个大型图像数据集的情况下，其中的图像已标记所含对象的类型、位置和大小，我们的目标是训练一个机器学习模型，该模型可以获取新的、未标记的图像，并检测对象是否存在，以及对象在图像中的位置和大小。

然而，概率模型中的隐变量并不需要有任何明确的物理解释，而引入这些隐变量可能只是为了从更简单的组件中构造出更复杂的联合分布。例如，归一化流、变分自编码器和扩散模型等都使用深度神经网络，从而通过转换具有简单高斯分布的隐变量，在数据空间中创建复杂的分布。稍后我们将看到扩散模型和其他生成式模型如何根据对所需图像内容和风格的文字描述，合成令人赞叹的高分辨率图像（参见第20章）。

\subsection{马尔可夫毯}\label{subsection:ch11-马尔可夫毯}

在讨论更复杂的有向图时，一个很有用的条件独立性称为马尔可夫毯（Markov blanket）或马尔可夫边界（Markov boundary）。考虑由具有$D$个节点的有向图表示的联合分布$p(x_1,\cdots,x_D)$，并考虑以所有其余变量$x_{j\neq i}$为条件的变量$x_i$的特定节点的条件分布。利用分解性质[式~\eqref{align:ch11-有向图联合分布}]，我们可以将这一条件分布表示为如下形式：

\begin{align}
  p(x_i\mid x_{j\neq i}) = \frac{p(\mathbf{x})}{\int p(\mathbf{x})\,\mathrm{d}x_i}
\end{align}

其中积分被离散变量的求和取代。现在我们注意到，任何与$x_i$无函数相关的因子$p(x_k\mid \mathrm{pa}(k))$都可以在对$x_i$的积分之外取值，因此分子和分母之间会抵消。剩下的因子就是节点$x_i$本身的条件分布$p(x_i\mid \mathrm{pa}(i))$，以及任何其他节点$x_k$的条件分布，从而节点$x_i$在$p(x_k\mid \mathrm{pa}(k))$的条件集中，即节点$x_i$是$x_k$的父节点。条件分布$p(x_i\mid \mathrm{pa}(i))$将取决于节点$x_i$的父节点，而条件分布$p(x_k\mid \mathrm{pa}(k))$将取决于$x_k$的子节点和共父节点[即节点$x_i$的父节点（节点$x_i$除外）所对应的变量]。由父节点、子节点和共父节点组成的节点集称为马尔可夫毯，如图11.25所示。它具有这样一个特性，即节点$x_i$的条件分布（以图中所有其余节点为条件）但取决于马尔可夫毯中的变量。我们可以考虑把节点$x_i$的马尔可夫毯看作能将节点$x_i$与图中的其他节点隔离开来的最小节点集。请注意，仅包含节点$x_i$的父节点和子节点是不够的，因为相消解释意味着观测到子节点不会阻塞通往共父节点的路径。因此，我们还必须观测到共父节点。

\subsection{作为过滤器的图}\label{subsection:ch11-作为过滤器的图}

我们已经看到，一个特定的有向图不仅代表了将联合概率分布分解为条件概率的乘积，还代表了通过d分离准则获得的一组条件独立性。d分离定理实际上表达了这两个特性的等价性。为了说明这一点，不妨把有向图看作一个过滤器。假设我们考虑的是图中（未观测到的）节点所对应的变量$\mathbf{x}$的特定联合概率分布$p(\mathbf{x})$。只有当且仅当该分布可以用图所隐含的分解[式~\eqref{align:ch11-有向图联合分布}]表示时，过滤器才会允许该分布通过。将变量$\mathbf{x}$上所有可能的分布$p(\mathbf{x})$的集合提交给过滤器，对于有向分解（directed factorization），将过滤器允许通过的分布子集记为$D_F$，如图11.26所示。我们也可以根据是否满足图的d分离特性时才能通过过滤器。将通过过滤器的所有可能概率分布$p(\mathbf{x})$的集合记为$D_F$。d分离特性所隐含的所有条件独立性，使用图（作为第二种过滤器）来过滤分布。d分离定理表明，同样的一组分布$D_F$将被允许通过第二种过滤器。如果我们将所有条件独立性，然后只允许满足所有这些条件独立性的分布通过。

需要强调的是，从d-分离得到的条件独立性适用于该特定有向图所描述的任何概率模型。例如，无论变量是离散的还是连续的，抑或是这些变量的组合，情况都是如此。在一种极端情况下，我们有一个全连接图，它完全不显示条件独立性，并且可以表示给定变量上任何可能的联合概率分布，集合$D_F$将包含所有可能的分布$p(\mathbf{x})$。在另一种极端情况下，我们有一个全不连接图，即不包含任何连接的图。这相当于联合分布是一个可以分解为组成图节点的变量的边缘分布的乘积。请注意，对于任何给定的图，集合$D_F$都将包括除图所描述的分布外，还具有其他独立特性的任何分布。例如，一个全分解分布总能通过相应变量集上的任何图所隐含的过滤器。

\end{document}
